\chapter{Alcance y delimitación del proyecto}

El alcance principal es el diseño y la evaluación experimental de un clasificador de lesiones cutáneas basado en aprendizaje autosupervisado. El clasificador es el módulo central del proyecto. La segmentación, el manejo de pelo/artefactos y la estimación de tono son módulos complementarios que se evalúan por separado; su integración con el clasificador requiere contratos de entrada/salida, gates aprobados y ablaciones que midan su aporte incremental.

HAM10000 constituye el conjunto etiquetado principal para la evaluación de clasificación de siete categorías: \texttt{nv}, \texttt{mel}, \texttt{bkl}, \texttt{bcc}, \texttt{akiec}, \texttt{vasc} y \texttt{df} \parencite{tschandl2018ham10000}. No contiene una clase independiente para carcinoma escamocelular invasor ni cubre patologías tropicales como leishmaniasis. El conjunto de preentrenamiento no supervisado y sus exclusiones se reportarán a partir de manifiestos y hashes de cada corrida. BCN20000 no se contará como un conjunto independiente cuando sus imágenes ya formen parte de ISIC2019.

PAD-UFES-20 se reporta como evaluación externa de dominio, sin ajuste del modelo, porque contiene imágenes clínicas capturadas con teléfonos y no es directamente comparable con la dermatoscopía de HAM10000 \parencite{pacheco2020padufes}. El informe describirá el protocolo, el mapeo de etiquetas y sus limitaciones. La pertenencia exacta de las imágenes a conjuntos, particiones y exclusiones se documentará con identificadores y hashes verificables.

El proyecto utiliza conjuntos públicos y no contempla la recolección de imágenes en hospitales ni la evaluación prospectiva con pacientes o profesionales clínicos. El prototipo se considera una demostración de laboratorio, no un dispositivo médico validado, una herramienta diagnóstica autónoma ni un sistema desplegado en una institución de salud. La validez externa, la representatividad regional y los análisis de subgrupos se limitarán a la evidencia observada y a la calidad de las anotaciones disponibles.




